\SourceLeaf{274}
\[
\begin{aligned}
\frac{\mathrm{d}U}{\mathrm{d}w}&=\frac{\cos\varphi}{\cos\theta}\\
\frac{\mathrm{dd}U}{\mathrm{d}w^2}&=-\frac{3ee\cos\varphi}{\cos\theta^3}\mathord{.}cs\\
\frac{\mathrm{d}^3U}{\mathrm{d}w^3}&=\frac{ee\cos\varphi}{\cos\theta^5}(3cc-3ss+ee(12ccss+3s^4))
\end{aligned}
\]
\[
\frac{\mathrm{d}^4U}{\mathrm{d}w^4}=\frac{ee\cos\varphi}{\cos\theta^7}\mathord{.}cs(16-ee(49cc-13ss)-e^4(56ccss+29s^4))
\]
\[
\begin{aligned}
\frac{\mathrm{d}^5U}{\mathrm{d}w^5}={}&\frac{ee\cos\varphi}{\cos\theta^9}(-16cc+12ss+ee(49c^4-378ccss+9s^4)\\
&+e^4(628c^4ss+174ccs^4-54s^6)+e^6(268c^4s^4+220ccs^6+33s^8))
\end{aligned}
\]
\indent Bei dieser Entwickelung von \(q\) in eine Reihe nach \(p\) ist stillschweigend vorausgesetzt, dass beide Grössen in Theilen des dem Halbmesser gleichen Bogens ausgedrückt sind: soll dagegen \(q\) Secunden und \(p\) Grade bedeuten, so muss dem ersten Gliede der Reihe der Factor \(3600\), dem zweiten der Factor \(\tfrac{3600\pi}{180}=20\pi\), dem dritten der Factor \(3600(\tfrac{\pi}{180})^2=\tfrac{1}{9}\pi\pi\) u.\,s.\,f. beigefügt werden. Unter dieser Voraussetzung gibt die Anwendung der Formeln auf unser Beispiel folgende Zahlenwerthe, welche ich in eine solche Form setze, dass weitgestreckte Decimalbrüche vermieden werden:
\[
\begin{aligned}
q={}&359556''69447\mathord{.}\frac{p}{100}\\
&+3041{,}386524\mathord{.}(\frac{p}{100})^2\\
&-946{,}260563\mathord{.}(\frac{p}{100})^3\\
&-4135{,}396057\mathord{.}(\frac{p}{100})^4\\
&+227{,}04342\mathord{.}(\frac{p}{100})^5
\end{aligned}
\]
welche Reihe, da \(p\) in der Anwendung nur wenige Einheiten betragen soll, immer sehr schnell convergirt. Um für die Richtigkeit dieser Zahlen eine Bestätigung zu erhalten, habe ich die Rechnung für \(p=-6\) und für \(p=+6\), d.\,i. für
\[
\begin{aligned}
P+p&=46^0\ 42'\ 2''53251\quad\text{und für}\\
P+p&=58\quad42\quad2{,}53251
\end{aligned}
\]
sowohl nach der Reihe, als nach der endlichen Formel 3 ausgeführt. Die Reihe gibt
